\subsection{Memorization}\label{subsec:memorization}

Numerous studies have demonstrated the feasibility of extracting training data from LLMs and diffusion models \citep{carlini2021extracting, carlini2023extracting, hu2022membership}. These models can memorize a substantial portion of their training data, raising concerns about potential breaches of confidential or copyrighted content. In the following section, we present initial exploration in the area of model memorization with RepE.

\subsubsection{Memorized Data Detection}
Can we use the neural activity of an LLM to classify whether it has memorized a piece of text? To investigate this, we conduct LAT scans under two distinct settings:
\begin{enumerate}[leftmargin=*]
    \item Popular vs. Synthetic Quotes: Popular quotes encompass well-known quotations sourced from the internet and human cultures. These quotes allow us to assess memorization of concise, high-impact text snippets. As a reference set, we prompt GPT-4 to generate synthetic quotations.
    \item Popular vs. Synthetic Literary Openings: Popular literary openings refer to the initial lines or passages from iconic books, plays, or poems. These openings allow us to assess memorization of longer text excerpts. As a reference set, we prompt GPT-4 to generate synthetic literary openings, modeled after the style and structure of known openings.
\end{enumerate}
Using these paired datasets as stimuli, we conduct LAT scans to discern directions within the model's representation space that signal memorization in the two settings separately. Since the experimental stimuli consist of likely memorized text which already elicits our target behavior, the LAT template does not include additional text. Upon evaluation using a held-out dataset, we observe that the directions identified by LAT exhibit high accuracy when categorizing popular and unpopular quotations or literary openings. To test the generalization of the memorization directions, we apply the directions acquired in one context to the other context. Notably, both of the directions transfer well to the other out-of-distribution context, demonstrating that these directions maintain a strong correlation with properties of memorization.


\begin{table}[t]

\centering
\begin{tabular}{lcccccccc}
\toprule
\multicolumn{1}{l}{} & \multicolumn{2}{c}{\multirow{2}{*} {No Control} } & \multicolumn{6}{c}{Representation Control}                                                             \\ \cmidrule(lr){4-9}
                      & \multicolumn{2}{c}{}                   & \multicolumn{2}{c}{Random}        & \multicolumn{2}{c}{$+$}           & \multicolumn{2}{c}{$-$} \\ \cmidrule(lr){2-3}\cmidrule(lr){4-5}\cmidrule(lr){6-7}\cmidrule(lr){8-9}
                      & EM                 & SIM                 & EM   & \multicolumn{1}{c}{SIM}   & EM   & \multicolumn{1}{c}{SIM}   & EM         & SIM        \\ \midrule
LAT\textsubscript{Quote} & \multirow{2}{*}{89.3} & \multirow{2}{*}{96.8} & 85.4 & \multicolumn{1}{c}{92.9} & 81.6 & \multicolumn{1}{c}{91.7} & 47.6 & 69.9 \\
LAT\textsubscript{Literature}            &                    &                    & 87.4 & \multicolumn{1}{c}{94.6} & 84.5 & \multicolumn{1}{c}{91.2} & \textbf{37.9}       & \textbf{69.8}      \\
\bottomrule
\end{tabular}
\caption{We demonstrate the effectiveness of using representation control to reduce memorized outputs from a LLaMA-2-13B model on the popular quote completion task. When controlling with a random vector or guiding in the memorization direction, the Exact Match (EM) rate and Embedding Similarity (SIM) do not change significantly. When controlled to decrease memorization, the similarity metrics drop significantly as the model regurgitate the popular quotes less frequently.}
\label{tab:mem_aa}
\end{table}


\subsubsection{Preventing Memorized Outputs}
Here, we explore whether the reading vectors identified above can be used for controlling memorization behavior. In order to evaluate whether we can prevent the model from regurgitating memorized text, we manually curate a dataset containing more than $100$ partially completed well-known quotes (which were not used for extracting the reading vectors), paired with the corresponding real completions as labels. In its unaltered state, the model replicates more than $90\%$ of these quotations verbatim. Following \Cref{subsec:rep-control}, we conduct control experiments by generating completions when applying the linear combination transformation with a negative coefficient using the reading vectors from the previous section. Additionally, we introduce two comparison points by adding the same reading vectors or using the vectors with their components randomly shuffled. The high Exact Match and Embedding Similarity scores presented in Table \ref{tab:mem_aa} indicate that using a random vector or adding the memorization direction has minimal impact on the model's tendency to repeat popular quotations. Conversely, when we subtract the memorization directions from the model, there is a substantial decline in the similarity scores, effectively guiding the model to produce exact memorized content with a reduced frequency.

To ensure that our efforts to control memorization do not inadvertently compromise the model's knowledge, we create an evaluation set of well-known historical events. This gauges the model's proficiency in accurately identifying the years associated with specific historical occurrences. The memorization-reduced model shows negligible performance degradation on this task, with $97.2\%$ accuracy before subtracting the memorization direction and $96.2\%$ accuracy afterwards. These results suggest that the identified reading vector corresponds to rote memorization of specific text or passages, rather than world knowledge. Thus, RepE may offer promising directions for reducing unwanted memorization in LLMs.





